We study multi-dealer secret sharing for threshold access structures. Such a scheme is parameterized by four variables:

\begin{itemize}
	\item $\mdssRecThr$, the number of shares required to recover a secret.
	\item $\mdssPrivThr$, the number of shares one can emit before privacy and unlinkability are broken.
	\item $\mdssMaxDealers$, the number of sufficient dealers the scheme is able to tolerate.
	\item  $\mdssMaxShares$, the maximum number of shares that can be input to the reconstruction algorithm.
\end{itemize}

We consider the \emph{ramp} setting of Blakley and Meadows~\cite{C:BlaMea84}, where we allow for a gap between $\mdssRecThr$ and $\mdssPrivThr$ larger than $1$.
We now formally define multi-dealer secret sharing.

\begin{definition}[Multi-Dealer Secret Sharing]
	A $(\mdssRecThr, \mdssPrivThr, \mdssMaxDealers, \mdssMaxShares)$-multi-dealer secre sharing scheme (MDSS) with security parameter $\secpar$ is parameterized by $(\mdssRecThr, \mdssPrivThr, \mdssMaxDealers, \mdssMaxShares)$, each of which is an integer polynomial in $\secpar$, and is defined over a secret space $\mdssSecretSpace$ and a share space $\mdssShareSpace$.
	\hme{Technically share and secret space should also be families over the security parameter}

	An MDSS scheme consists of the following algorithms, which are expected to run in time polynomial in $\secpar$.

	\begin{itemize}
		\item $\mdssShareAlg(\secparam, \mdssSecret, \mdssNumShares) \rightarrow \set{\mdssShare{\mdssSecret}_i}_{i \in [\mdssNumShares]}$, takes as input a secret $\mdssSecret \in \mdssSecretSpace$ and an integer $\mdssNumShares$ and outputs a set of shares $\set{\mdssShare{\mdssSecret}_i}$ where each $\mdssShare{\mdssSecret}_i \in \mdssShareSpace$.
		\item $\mdssReconstructAlg(\mdssShare{\mdssSecret}_1, \dots, \mdssShare{\mdssSecret}_{\mdssWindowSize}) =: \set{\mdssShare{\mdssSecret}_i}$, takes as input shares $\mdssShare{\mdssSecret}_1$, \dots, $\mdssShare{\mdssSecret}_{\mdssWindowSize}$ where $\mdssWindowSize \leq \mdssMaxShares$ and outputs a (potentially empty) set of secrets.
	\end{itemize}
\end{definition}

We now define some notational shorthand.
For a secret $\mdssSecret$, an integer $\mdssNumShares$, and a list of unique indices $\indexSet \subseteq [\mdssNumShares]$, define the algorithm $\mdssProj(\mdssSecret, \mdssNumShares, \indexSet)$ as follows:
\procedureblock{$\mdssProj(\mdssSecret, \mdssNumShares, \indexSet)$}{
	\set{\mdssShare{\mdssSecret}_1, \dots, \mdssShare{\mdssSecret}_{\mdssNumShares}} \gets \mdssShareAlg(\mdssSecret, \mdssNumShares) \\
	\pcreturn \set{\mdssShare{\mdssSecret}_i}_{i \in \indexSet}
}

Next, for sets of shares $\mdssShareSet_1, \dots, \mdssShareSet_\mdssNumShareSets$, define the algorithm $\mdssFlatten(\mdssShareSet_1, \dots, \mdssShareSet_\mdssNumShareSets)$ as the function that returns a combined list of the shares in each share set.
For example, $\mdssFlatten\left(\set{\mdssBlindShare^1_1, \mdssBlindShare^1_2}, \set{\mdssBlindShare^2_1}\right) = \mdssBlindShare^1_1, \mdssBlindShare^1_2, \mdssBlindShare^2_1$.

Finally, define $\mdssGenerate(\set{(\mdssSecret_i, \mdssNumShares_i, \indexSet_i)}_{i \in [\mdssNumShareSets]})$ as the algorithm that runs $\mdssProj$ on each configuration in its input, then returns the flattened set of shares.

\begin{definition}[Unlinkability]
	Let $(\mdss = (\mdssShareAlg, \mdssReconstructAlg)$ be an MDSS scheme and let the random vaiable $\ulink_b(\mdss, \adv, \secpar)$ where $\adv = \set{\adv_1, \adv_2}$, $b \in \bin$, and $\secpar \in \NN$ denote the result of the following experiment:
	\procedureblock{$\ulink_b(\pke, \adv, \secpar)$}{
		\set{(\mdssSecret_i, (\mdssNumShares_{i, 0}, \indexSet_{i, 0}), (\mdssNumShares_{i, 1}, \indexSet_{i, 1}))}_{i \in [\mdssNumShareSets]}, \st \gets \adv_1(\secparam) \\
		\mdssBlindShare_1,\dots,\mdssBlindShare_{\mdssWindowSize} \gets \mdssGenerate(\set{(\mdssSecret_{i,b}, \mdssNumShares_{i, b}, \indexSet_{i, b})}_{i \in [\mdssNumShareSets]}) \\
		b' \gets \adv_2(\st, \mdssBlindShare_1, \dots, \mdssBlindShare_\mdssWindowSize) \\
		\pcreturn b'
	}
	We say that an adversary $\adv$ is \emph{admissible} if it satitfies:
	\begin{enumerate}
		\item $\forall i \in [\mdssNumShareSets], \forall b \in \bin, |\indexSet_{i, b}| \leq \mdssPrivThr$
		\item $\sum_{i \in [\mdssNumShareSets]} |\indexSet_{i, 0}| = \sum_{i \in [\mdssNumShareSets]} |\indexSet_{i, 1}|$
	\end{enumerate}

	The $\ulink(\mdss, \adv, \secpar)$ advantage of $\adv$ is defined as
	$$
		\advantage{\ulink}{\adv, \mdss} = \abs{\prob{\ulink_0(\pke, \adv, \secpar) = 1} - \prob{\ulink_1(\pke, \adv, \secpar) = 1}}
	$$
	We say that an MDSS scheme is $\mdssULVar$-unlinkable if for all admissible NUPPT adversaries $\adv$ we have that $\advantage{\ulink}{\adv, \mdss} \leq \mdssULVar(\secpar)$.
\end{definition}


While the above unlinkability definition is powerful, it is also somewhat cumbersome.
We therefore define a simpler definition for which it will be easier to prove that a scheme satisfies, and then show that the two definitions are in fact equivalent.

\begin{definition}[One-More Unlinkability]
	We say that an MDSS scheme is $\mdssULVar$\emph{-one-more-unlinkable} if for all admissible NUPPT adversaries $\adv = (\adv_1, \adv_2)$ and for all $\secpar \in \NN$:
	\begin{equation*}
		\begin{gathered}
			\probsublong{
				\begin{gathered}
					(\mdssSecret_0, \mdssSecret_1, \mdssNumShares_0, \mdssNumShares_1, \indexSet, i), \st \gets \adv_1(\secparam) \\
					\set{\mdssShare{\mdssSecret_0}_1, \dots, \mdssShare{\mdssSecret_0}_{\mdssNumShares_0 - 1}, \mdssShare{\mdssSecret_0}^*} \gets \mdssProj(\mdssSecret_0, \mdssNumShares_0, \indexSet)\\
					\set{\mdssShare{\mdssSecret_1}^*} \gets \mdssProj(\mdssSecret_1, \mdssNumShares_1, i) \\
				\end{gathered}
			}{
				\adv_2(\st, \mdssShare{\mdssSecret_0}_1, \dots, \mdssShare{\mdssSecret_0}_{\mdssNumShares_0 - 1}, \mdssShare{\mdssSecret_b}^*) = b
			} \\
			\leq \frac{1}{2} + \mdssULVar(\secpar)
		\end{gathered}
	\end{equation*}
\end{definition}
where $\adv$ is \emph{admissible} if it is the case that $|\indexSet| \leq \mdssPrivThr$.
We refer to $\mdssULVar(\secpar)$ as the \emph{advantage} of $\adv$, and denote it by $\mdssOMUAdv$.

As with Unlinkability we can use the standard transformation to produce a two-words definition, use the notation $\mdssOMULGame{\hat{b}}{\adv}$, and have that $\mdssOMUAdv = \left| \prob{\mdssOMULGame{0}{\adv} = 1} - \prob{\mdssOMULGame{1}{\adv} = 1}\right|$.

\begin{definition}[MD-Correctness]
	We say that a $(\mdssRecThr, \mdssPrivThr, \mdssMaxDealers, \mdssMaxShares)$-MDSS scheme is $\mdssCVar$\emph{-correct} if for all $\secpar \in \NN$ and for all $\set{(\mdssSecret_i, \mdssNumShares_i, \indexSet_i)}_{[\mdssNumShareSets]}$ where the following hold:
	\begin{itemize}
		\item $\forall i \in [\mdssNumShareSets],\: \indexSet_i \subseteq [\mdssNumShares_i]$

		\item $\sum_{i \in [\mdssNumShareSets]} |\indexSet_i| \leq \mdssMaxShares$
		\item $|\set{\indexSet_i \text{ s.t. } |\indexSet_i| \geq \mdssRecThr}| \leq \mdssMaxDealers$

		\item $\bigcap_{i, j \in [\mdssNumShareSets],\: i \neq j} \indexSet_i \cap \indexSet_j = \emptyset$
	\end{itemize}
	it holds that:
	\begin{equation*}
		\probsublong{
			\mdssShareList \gets  \mdssGenerate\left(\set{(\mdssSecret_i, \mdssNumShares_i, \indexSet_i)}_{i \in [\mdssNumShareSets]}\right)
		}{
			\mdssReconstructAlg(\mdssShareList) \neq \set{\mdssSecret_i \text{ s.t. } |\indexSet_i| \geq \mdssRecThr}
		} \leq \mdssCVar(\secpar)
	\end{equation*}
\end{definition}
